\documentclass[11pt]{article}
\usepackage[margin=1in]{geometry}
\usepackage[T1]{fontenc}
\usepackage{lmodern,amsmath,amssymb,amsthm,mathtools,booktabs,longtable,array}
\usepackage{microtype,enumitem,xcolor,fancyhdr}
\usepackage[colorlinks=true,linkcolor=blue!45!black,citecolor=blue!45!black,urlcolor=blue!45!black]{hyperref}
\hypersetup{pdftitle={Fewer integer variables in Diophantine formulations of ZFC consistency},pdfauthor={OpenAI-assisted research synthesis for ProveIt}}
\pagestyle{fancy}\fancyhf{}\fancyhead[L]{\small Diophantine consistency: 27 integer variables}\fancyhead[R]{\small ProveIt research note}\fancyfoot[C]{\thepage}
\setlength{\headheight}{14pt}
\setlength{\emergencystretch}{2em}
\setcounter{tocdepth}{1}
\newtheorem{theorem}{Theorem}[section]
\newtheorem{lemma}[theorem]{Lemma}
\newtheorem{proposition}[theorem]{Proposition}
\newtheorem{corollary}[theorem]{Corollary}
\theoremstyle{definition}\newtheorem{definition}[theorem]{Definition}
\theoremstyle{remark}\newtheorem{remark}[theorem]{Remark}
\newcommand{\Z}{\mathbb Z}\newcommand{\N}{\mathbb N}
\newcommand{\Pos}{\operatorname{Pos}}\newcommand{\Con}{\operatorname{Con}}
\newcommand{\Res}{\mathcal R}\newcommand{\pc}{\operatorname{pc}}
\newcommand{\sg}{\operatorname{sgn}}
\newcommand{\code}[1]{\texttt{\detokenize{#1}}}
\title{Fewer integer variables in Diophantine\ formulations of ZFC consistency\\[5pt]\large Weighted root amalgamation, 27 unrestricted integers,\ and a smaller finite compiler alphabet}
\author{OpenAI-assisted research synthesis for the ProveIt repository}
\date{22 October 2026}
\begin{document}
\maketitle
\begin{abstract}
We combine and sharpen three constructions in ProveIt's current collection of explicit Diophantine encodings of ZFC consistency. First, a new weight in the three-root resultant preserves 13 positive unknowns and 84 binary arithmetic operations while lowering the degree from 571 to 515. The same degree previously required 85 operations. Second, a signed reconstruction theorem combines that resultant with the seven-positive-coordinate reduction: an even weight and an odd primitive element give a polynomial in 27 unrestricted integers, improving the collection's previous 29. Its exact degree is 1004 and its circuit has 139 operations; a four-equation version has maximum degree 952 and 133 operations. Third, a local pruning of the four-cell compiler reduces its alphabet from 37,075 to 5,268 tiles and its fixed radix from $2^{5^{18}}$ to $2^{5^{17}}$. We give an exact finite coefficient recipe, ordinary proofs of the new reductions, executable circuits, and finite verification artifacts. For the resulting fixed polynomials, ZFC proves that absence of a solution is equivalent to $\Con(\mathrm{ZFC})$, by composition with the cited compiler and machine-calibration theorems. No global minimality, consistency proof, or new end-to-end proof-assistant verification is claimed.
\end{abstract}
\begingroup\small\tableofcontents\endgroup

\section{Result, comparison, and scope}
\label{sec:scope}
The source snapshot is the ProveIt commit
\code{f0ae4e0991ca44a7bb4e5671d42e33d775e48035}.
Its reconciliation note records seven overlapping packages, numbered 09--15, with different optimization objectives~\cite{collection}. The existing best variable count is 13 for positive integers and 29 for unrestricted integers. Package 15 is not the frontier in every metric. We compare against the whole collection, particularly packages 09, 11, and 14.

\begin{theorem}[Explicit constructions]
\label{thm:main}
The coefficient tuple specified in Section~\ref{sec:coeff} determines ordinary polynomials
\[
P_{13}\in\Z[x_1,\ldots,x_{13}],\qquad
P_{27}\in\Z[z_1,\ldots,z_{27}]
\]
such that
\begin{align}
\mathrm{ZFC}&\vdash
 \bigl(\neg\exists\mathbf x\in\Z_{>0}^{13}\ P_{13}(\mathbf x)=0\bigr)
 \longleftrightarrow\Con(\mathrm{ZFC}),\label{eq:theorem13}\\
\mathrm{ZFC}&\vdash
 \bigl(\neg\exists\mathbf z\in\Z^{27}\ P_{27}(\mathbf z)=0\bigr)
 \longleftrightarrow\Con(\mathrm{ZFC}).\label{eq:theorem27}
\end{align}
Their exact total degrees are 515 and 1004, respectively. Their division-free straight-line programs have 84 and 139 binary operations, respectively, counting multiplication by fixed numerals. Each also has an equivalent four-equation existence formulation, with the dimensions in Table~\ref{tab:frontier}.
\end{theorem}

The proof of this theorem consists of the new algebra in Sections~\ref{sec:polynomial}--\ref{sec:signed}, the new finite-language adapter in Section~\ref{sec:tiles}, and the established interfaces stated in Section~\ref{sec:inherited}. These are mathematical proofs formalizable in ZFC. The archive does not contain a new Lean or Rocq proof of the Diophantine theorem.

\begin{table}[ht]
\centering\small
\begin{tabular}{@{}lrrrr@{}}
\toprule
Construction & Unknowns & Equations & Degree & Operations\\
\midrule
09: original weighted resultant & $13>0$ & 1 & 571 & 84\\
09: previous smaller weight & $13>0$ & 1 & 515 & 85\\
\textbf{This note: smaller weight} & $\mathbf{13>0}$ & 1 & \textbf{515} & \textbf{84}\\
09: original resultant system & $13>0$ & 4 & 544 & 78\\
\textbf{This note: positive system} & $\mathbf{13>0}$ & 4 & \textbf{488} & \textbf{78}\\
14: ordinary integer product & $29\Z$ & 1 & 188 & 121\\
\textbf{This note: integer product} & $\mathbf{27\Z}$ & 1 & \textbf{1004} & \textbf{139}\\
\textbf{This note: integer system} & $\mathbf{27\Z}$ & 4 & \textbf{952} & \textbf{133}\\
\bottomrule
\end{tabular}
\caption{Comparison at the pinned snapshot. A system's degree is the maximum residual degree. The reduction from 29 to 27 integer unknowns trades degree and arithmetic size for fewer variables.}
\label{tab:frontier}
\end{table}

This does not dominate every prior result. The 15-positive-variable, 67-operation baseline remains shorter as an evaluator. The existing 32-positive-variable quartic remains much lower in degree. The present result reduces the unrestricted-integer variable count and improves the joint degree/operation profile at 13 positive variables. It also reduces the fixed coefficients. The repository's existential nine-unknown theorem is a different result: it does not supply the specialized, fully specified circuit being compared here~\cite{collection}.

\section{The inherited arithmetic and logical interfaces}
\label{sec:inherited}
The purpose of an explicit interface is to show exactly what is reused. We do not infer an internal ZFC theorem merely from the existence of an external Lean certificate.

\subsection{Exact machine endpoint}
Let $K_{215}$ be the literal blank-input two-symbol machine in the archive. Its 430 transitions have SHA-256
\begin{center}\small\code{fdbcef0825886035c3091c8537a69535bc96f9e265f0ae45863543d40ae70090}.\end{center}
State 0 is initial, symbol 0 is blank, and state 215 is halting. Its actual initial action writes 1, moves right, and enters state 93. The repository's mathematical calibration proves in ZFC
\begin{equation}
\operatorname{Halt}(K_{215},\mathrm{blank})
 \longleftrightarrow\neg\Con(\mathrm{ZFC}).\label{eq:machine}
\end{equation}
The Lean endpoint independently has the exact type
\begin{center}\small
\code{BlankHalts machine215} $\longleftrightarrow$ \code{BProv ZFCax_s [] Form.fBot}.
\end{center}
The semantic chain and its ordinary formalizability are described in the calibration and verification documents~\cite{machine,endpoint,verification}. No consistency or model-existence hypothesis is assumed in this biconditional. External kernel checking is valuable evidence, but is not itself a derivation encoded in ZFC's proof syntax.

\subsection{Four fixed coefficient ports}
The arithmetic uses one complete tuple
\[
(\beta,K_0,M_C,M_F^*),\qquad
\beta=B-1,\quad K_0=D_C+BD_R,\quad M_F^*=M_F+B-1.
\]
Here $M_F$ is the native mask; the arithmetic port is its \emph{shifted} version $M_F^*$. All four entries must come from the same compiler recipe. The native masks satisfy
\begin{equation}
0<M_C,M_F<B-1,\quad M_C\equiv2\pmod4,\quad M_F\equiv4\pmod8,
\quad\pc(M_C)+\pc(M_F)=\log_2B.\label{eq:masks}
\end{equation}
Section~\ref{sec:coeff} gives a new authentic tuple, using the same compiler contract as packages 06, 09, 11, and 14.

\subsection{Normalized Pell system and the signed-domain theorem}
Use the abbreviations to be defined in \eqref{eq:common}. In the inherited normalized system, three additional roots $f,\tau,y$ satisfy
\begin{align}
\tau^2&=A,\label{eq:oldfirst}\\
f^2&=1+\Delta i^2c^4,\label{eq:oldstrong}\\
Q V^2-(Q-1)y^2&=1,\qquad
 Q=\Delta^2i^2c^4,\quad V=c(Tf-1)-Rf^2.\label{eq:oldaux}
\end{align}
The remaining factors are $N_1,N_k,N_t$. The exact signed extension established by package 14 permits
\[
J,F,\alpha,z,s,\eta,\zeta>0,\qquad i\text{ odd},\quad w\text{ a square},
\]
and unrestricted signs for $f,h,T,\tau,y,\gamma$. Its product form is
\begin{equation}
N_0N_1N_aN_kN_tN_s=1,\label{eq:signedcontract}
\end{equation}
where
\[
N_0=\tau^2-(A-1),\quad N_s=f^2-\Delta i^2c^4,\quad
N_a=QV^2-(Q-1)y^2.
\]
Existence of a tuple in this domain is equivalent to blank halting of $K_{215}$~\cite{integer29,blank,pell}.

For clarity, the soundness theorem includes signed auxiliary quotients and signed main-root offsets; simply discarding positivity in an earlier positive-variable theorem would not suffice. Its proof uses normalized strong divisibility, signed Pell index congruences, the lower Pell-ratio estimate, and exact population inversion. After reconstruction it forces $N_k=N_t=1$ and the correct computation index. The complete cited signed proof is bundled as \code{upstream/SIGNED_DOMAIN_PROOF.md}.

The canonical completeness witnesses can be chosen with $R\equiv3\pmod4$, $i$ odd, $w$ a square, and $y$ odd. The transfer of these properties to the new coefficient tuple is proved in Section~\ref{sec:parity}. These are choices of witnesses, not assertions that all original positive solutions have those properties.

\section{The explicit polynomial}
\label{sec:polynomial}
All names introduced below are polynomial abbreviations. They are substituted expressions, not extra quantified unknowns or unevaluated predicates.

\subsection{Common arithmetic}
For coordinates $J,F,\alpha,z,h,i,T,s,w,\eta,\zeta,\gamma,t$, define
\begin{gather}
\begin{aligned}
q&=\beta J+1,&\qquad C_w&=q-F-\alpha,\\
X&=wq,& Y&=sq^3,\\
E&=XY,& k&=\eta+\zeta,\\
c&=kY+\eta,& a&=E+Y,\\
H&=4a+3,& \Delta&=a^2+H,\\
A_0&=a+2,& D_m&=X+ac+\gamma H,
\end{aligned}\notag\\
R=((q-1)(q-F)+\alpha)(q^2-1)+(M_C+qM_F^*)J,\notag\\
N_1=D_m^2-\Delta c^2,\qquad N_k=k-hE-R,\notag\\
N_t=(K_0+w)C_w+q-F-z(q-1).
\label{eq:common}
\end{gather}
In particular $\Delta=A_0^2-1$. Next define
\begin{align}
\ell&=EkY,& A&=1+\ell(\ell+k),& I&=ic^2,\notag\\
D&=1+\Delta I^2,& Q&=\Delta(D-1)=\Delta^2 I^2,& d&=Q-1,\notag\\
v&=cT,& z_0&=c+RD,&
P_0&=QDv^2,& Z_0&=Qz_0^2,\notag\\
e&=P_0+Z_0-1.& &&
\label{eq:rootdata}
\end{align}
For a weight $M$ specified below, set
\begin{align}
L&=dt^2,& N&=dM^2A,& C&=L+N-e,\notag\\
U&=C^2+4(P_0Z_0-LN),&
\Res&=U^2-16C^2P_0Z_0.\label{eq:resultant}
\end{align}
The common finalizer is
\begin{equation}
\boxed{\ P=N_1N_kN_t(1+\Res)-1.\ }\label{eq:polynomial}
\end{equation}

\subsection{Thirteen positive integers}
For $P_{13}$, all thirteen listed coordinates range over $\Z_{>0}$, and
\begin{equation}
\boxed{M=I^2+v^2.}\label{eq:positiveweight}
\end{equation}
Equations \eqref{eq:common}--\eqref{eq:positiveweight} are a complete polynomial recipe once the four fixed integers are specified. In particular there is no variable exponentiation, divisibility predicate, unbounded product, or implicit search during polynomial evaluation.

\subsection{Twenty-seven unrestricted integers}
Define the positive surjection
\begin{equation}
\Pos(x_0,x_1,x_2)=x_0^2+x_0+x_1^2+x_2^2+1.\label{eq:pos}
\end{equation}
Use seven disjoint copies of \eqref{eq:pos} to define
\begin{equation}
J,F,\alpha,z,s,\eta,\zeta\label{eq:seven}
\end{equation}
Retain six independent unrestricted coordinates
\[
h,r_i,T,u,\gamma,r_t,
\]
and substitute
\begin{equation}
i=2r_i+1,\qquad w=u^2,\qquad t=2r_t+1,\qquad
\boxed{M=2(I^2+v^2).}\label{eq:integerweight}
\end{equation}
This defines $P_{27}$ literally by substitution in \eqref{eq:polynomial}. Its number of unrestricted integer variables is
\[
7\cdot3+6=27.
\]
The even weight and the odd primitive element are essential parts of this domain conversion: oddness makes $t$ nonzero, while the even weight makes canonical completeness compatible with that restriction.

\begin{lemma}[Positivity surjection]
The map \eqref{eq:pos} has image exactly $\Z_{>0}$.
\end{lemma}
\begin{proof}
For every integer $x_0$, $x_0(x_0+1)\ge0$, so its value is positive. For $n\ge1$, the number $4n-3$ is positive and is either 1 or 5 modulo 8. Legendre's three-square theorem applies. Modulo 4 exactly one representing square is odd. Write the three squares as $(2x_0+1)^2,(2x_1)^2,(2x_2)^2$ and divide the resulting identity by 4.
\end{proof}

\section{The new denominator bound}
\label{sec:denominator}
A resultant equation alone can have spurious solutions when a reconstruction denominator vanishes. The central new observation is that the retained main Pell norm supplies enough growth to use the smaller weight $I^2+v^2$.

\subsection{Units and the weak transport bootstrap}
At a zero of \eqref{eq:polynomial}, each of its four integer factors is a unit. Since $\Delta\equiv0$ or $3\pmod4$, the factor $N_1$ cannot equal $-1$. Since $\Res$ is a square modulo 4, $1+\Res$ cannot equal $-1$. Thus
\begin{equation}
N_1=1,\qquad\Res=0,\qquad N_k=N_t=\epsilon\in\{1,-1\}.\label{eq:units}
\end{equation}

For the integer construction, initially $w\ge0$ and $Y,k,c,a>0$. If $w=0$, then $X=E=0$, $a=Y$, and $D_m=ac+\gamma H$. Reducing $D_m^2-\Delta c^2=1$ modulo $H$ gives $0\equiv1\pmod H$, impossible because $H=4Y+3>1$. Hence
\begin{equation}
w>0,\qquad X\ge q\ge16.\label{eq:zerow}
\end{equation}
This also holds automatically in the positive construction.

Weak transport gives only $C_w\ge0$ at this point. Indeed
\[
q-F-z(q-1)=1-F-(z-1)(q-1)\le0,
\]
so $C_w\le-1$ would give $N_t\le-(K_0+w)\le-2$. Thus $q-F=C_w+\alpha>0$. With
\[
G=(q-1)(q-F)+\alpha=q\alpha+(q-1)C_w,
\]
the actual mask bounds imply
\[
q\le G\le q^2-q,\qquad
0<(M_C+qM_F^*)J<2q^2-q-1.
\]
Consequently the bootstrap needed below is
\begin{equation}
3q+1<R<q^4-2,\quad E\ge q^4>R+2,\quad a>R+2,
\quad Y\ge q^3\ge4096.\label{eq:bootstrap}
\end{equation}
No deleted root or decoded computation is used here. In particular, $C_w=0$ is not prematurely excluded.

\subsection{The main Pell index cannot be two}
For $b\ge2$, define the integer Pell sequences by
\[
\chi_b(n)+\psi_b(n)\sqrt{b^2-1}
 =(b+\sqrt{b^2-1})^n.
\]
The standard classification of the positive integer solutions of the Pell equation for $b^2-1$ gives an index $p$ with
\[
c=\psi_{A_0}(p),\qquad |D_m|=\chi_{A_0}(p).
\]
Because $c>1$, $p\ge2$.

\begin{lemma}[Early index growth]
\label{lem:early}
Under \eqref{eq:units}--\eqref{eq:bootstrap}, even if $\gamma$ is signed, $p\ge3$. Therefore
\begin{equation}
c\ge4A_0^2-1>3A_0^2.\label{eq:cgrowth}
\end{equation}
\end{lemma}
\begin{proof}
Suppose $p=2$. Then $c=2A_0$ and $D_0=\chi_{A_0}(2)=2A_0^2-1$, with
\[
D_0-ac=H+4.
\]
If $D_m=D_0$, its defining congruence gives $X\equiv4\pmod H$, impossible since $16\le X<a<H$. If $D_m=-D_0$, use
\[
(D_0-ac)(D_0+ac)=1+Hc^2.
\]
It gives $4(D_0+ac)\equiv1\pmod H$ and hence $4X\equiv-1\pmod H$. But $0<4X+1<H$ because $a=(X+1)Y>X$. This also is impossible. Now $\psi_{A_0}(3)=4A_0^2-1$ and the Pell sequence increases.
\end{proof}

\begin{lemma}[Small-weight domination]
\label{lem:domination}
Let $i$ be nonzero, allow any integer $T$, and put $M_0=I^2+v^2$. Under the preceding hypotheses,
\begin{equation}
AM_0^2>2(Dv^2+z_0^2).\label{eq:domination}
\end{equation}
For either $M=M_0$ or $M=2M_0$, every integer $t$ therefore satisfies $C>0$.
\end{lemma}
\begin{proof}
The positive variables $\eta,\zeta$ give $kY<c<k(Y+1)$. By Lemma~\ref{lem:early}, $k>3A_0^2/(Y+1)$. Since $\ell=kXY^2$,
\begin{align*}
A&>k^2X^2Y^4\\
 &>9A_0^4E^2\left(\frac{Y}{Y+1}\right)^2
 \ge\frac{144}{25}A_0^4E^2
 >4(R+1)^2A_0^4.
\end{align*}
Here $Y\ge4$ and $E>R+2$. Also $I^2\ge c^4>c>1$ and
\[
D=1+(A_0^2-1)I^2<A_0^2I^2,
\qquad z_0<(R+1)A_0^2I^2.
\]
The inequalities $M_0^2\ge I^4$ and $M_0^2\ge4I^2v^2$ now yield
\[
AM_0^2>4z_0^2,\qquad AM_0^2>4Dv^2.
\]
For the second inequality use $A>A_0^2$; when $v=0$ it is immediate. Averaging proves \eqref{eq:domination}. Finally $Q>1$ and $d=Q-1\ge Q/2$, so
\[
dAM^2>Q(Dv^2+z_0^2)=e+1,
\qquad C=dt^2+dAM^2-e>dt^2+1>0.
\]
\end{proof}

\section{Exact reconstruction and the 27-integer theorem}
\label{sec:signed}
The next proof treats the signed extension, including a denominator that is automatically nonzero in the positive case.

\subsection{Excluding the zero auxiliary quotient}
\begin{lemma}
\label{lem:nonzeroT}
For a zero of $P_{27}$, $T\ne0$.
\end{lemma}
\begin{proof}
We know $d>0$, $M>0$, $t$ odd and nonzero, and $C>0$. First note
\begin{align*}
Q-z_0&=(\Delta-R)D-\Delta-c\\
&\ge D-\Delta-c
 =1+\Delta(I^2-1)-c>1.
\end{align*}
Thus $1<z_0<Q-1$; here $\Delta>R$ follows already from \eqref{eq:bootstrap}. Suppose $T=0$. Then $v=P_0=0$, and $\Res=U^2=0$ gives $U=C^2-4LN=0$. The rational number
\[
\tau=\frac{C}{2dtM}
\]
has square $A$. A rational number whose square is an integer is itself an integer: in lowest terms its denominator must divide its numerator. Set $y=t-M\tau\in\Z$. Expanding the defining relation gives
\[
Qz_0^2-(Q-1)y^2=1.
\]
It follows that $(Q-1)(y^2-z_0^2)=z_0^2-1>0$, so $|y|\ge z_0+1$ and
\[
(Q-1)(2z_0+1)\le z_0^2-1.
\]
This contradicts $Q-1>z_0>1$.
\end{proof}

\subsection{Recovering the three deleted roots}
Introduce the proof-only expression $B_1=2Qvz_0$. The polynomial identities
\begin{align}
DB_1^2&=4P_0Z_0,\notag\\
U&=C^2+DB_1^2-4d^2t^2M^2A,\notag\\
\Res&=U^2-D(2CB_1)^2\label{eq:rootidentity}
\end{align}
follow directly from \eqref{eq:rootdata}--\eqref{eq:resultant}.
All denominators in the following formulas are now established nonzero:
\begin{equation}
f=-\frac{U}{2CB_1},\qquad
\tau=\frac{C+B_1f}{2dtM},\qquad
y=t-M\tau.\label{eq:inverse}
\end{equation}
From $\Res=0$ we get $f^2=D$, and rational-square integrality gives $f\in\Z$. Using $U+2CB_1f=0$ gives
\[
(C+B_1f)^2=4d^2t^2M^2A,
\]
so $\tau^2=A$ and $\tau,y\in\Z$. Finally
\begin{align}
dy^2&=e-B_1f=Q(vf-z_0)^2-1,\notag\\
f^2&=1+\Delta i^2c^4,\notag\\
vf-z_0&=c(Tf-1)-Rf^2.\label{eq:restored}
\end{align}
These are precisely \eqref{eq:oldfirst}--\eqref{eq:oldaux}. The retained factors satisfy $N_1=1$ and $N_kN_t=1$, so the restored tuple satisfies the inherited signed product \eqref{eq:signedcontract}. Its soundness theorem gives blank halting of $K_{215}$.

\subsection{Completeness and parity}
If $K_{215}$ halts, choose a canonical witness with the properties in Section~\ref{sec:parity}. In particular $y$ is odd, $i$ is odd, and $w$ is a square. Evaluate $M=2(I^2+v^2)$ and put
\[
t=M\tau+y.
\]
It is odd and thus has the form $2r_t+1$. The identities behind \eqref{eq:inverse} give $\Res=0$. The other factors all equal 1. The positive surjection realizes the seven coordinates, so $P_{27}$ has an unrestricted-integer zero.

Together with reconstruction, this proves
\[
\exists\mathbf z\in\Z^{27}\ P_{27}(\mathbf z)=0
\quad\longleftrightarrow\quad\operatorname{Halt}(K_{215},\mathrm{blank}).
\]
Composition with \eqref{eq:machine} and classical negation yields \eqref{eq:theorem27} internally in ZFC.

\subsection{The positive construction and the four equations}
For $P_{13}$ use the same proof with $M=M_0$ and positive $t,i,T,w$. Lemma~\ref{lem:nonzeroT} is unnecessary because $T>0$. Reconstruction may return a negative $f$; replacing $f,T$ simultaneously by their negatives preserves $f^2$ and $Tf$. Replace $\tau,y$ by their positive absolute values: their equations use only squares, while $A>1$ and the auxiliary norm with $Q>1$ exclude zero. The signed auxiliary-quotient soundness theorem used in package 09 now applies~\cite{amalgamation,signedquotient}. The new weight changes only the denominator bound in that reconstruction; the other retained hypotheses are unchanged. Thus no positivity of a reconstructed quotient is silently assumed.

For completeness, positive canonical roots give $t=M_0\tau+y>0$. This proves \eqref{eq:theorem13}. The four equations
\begin{equation}
\boxed{N_1=1,\qquad N_k=1,\qquad N_t=1,\qquad\Res=0}\label{eq:system}
\end{equation}
have the same existence meaning in each stated domain. Indeed they imply $P=0$, while halting gives canonical witnesses satisfying all four. The decoded product argument also forces the two remaining unit signs positive. The circuit files represent the last equation as $U^2=16C^2P_0Z_0$, saving its final subtraction.

\section{Exact arithmetic sizes}
\label{sec:degree}
The counts refer to a straight-line program with reuse of previous registers. Every binary addition, subtraction, and multiplication is charged, including multiplication by 2, 4, or 16. Loading a fixed integer and using an already computed register are free. A system's comparisons are not counted as arithmetic gates; its residual degrees are still reported.

\subsection{Why the new weight saves a gate at degree 515}
Compute the strong registers as
\[
I=ic^2,\quad I_2=I^2,\quad D_-=\Delta I_2,
\quad D=D_-+1,\quad Q=\Delta D_-.
\]
These five gates replace five gates in package 09. The weight $M_0=I_2+v^2$ uses one addition, just as the old $D+v^2$ did. Thus the positive system keeps 78 gates and the product keeps 84. Package 09's earlier degree-515 weight $kI^2+T^2$ used one additional gate. The changed weight is justified by Lemma~\ref{lem:early}, which was not required for the previous larger weight.

\begin{table}[ht]
\centering\small
\begin{tabular}{@{}lrr@{}}
\toprule
Expression & Positive degree & Integer-substitution degree\\
\midrule
$q,X,Y,E,k,c,\Delta$ & $1,2,4,6,1,5,12$ & $2,4,8,12,2,10,24$\\
$R,I^2,A,D,Q,v,z_0$ & $4,22,22,34,46,6,38$ & $8,42,44,66,90,11,74$\\
$P_0,Z_0,M,L,N,C$ & $92,122,22,48,112,122$ & $178,238,42,92,218,238$\\
$P_0Z_0,LN,U,\Res$ & $214,160,244,488$ & $416,310,476,952$\\
$N_1,N_k,N_t$ & $18,7,2$ & $35,13,4$\\
$P$ & $515$ & $1004$\\
\bottomrule
\end{tabular}
\caption{Uniform exact degrees for every authentic fixed tuple, with $\beta>0$. Doubling the integer weight changes no degrees.}
\label{tab:degrees}
\end{table}

These are mathematical exact-degree claims, not only numerical diagnostics. In both columns, $C$ has unique highest part $-(Z_0)_{\rm top}$, $U$ has unique highest part $C_{\rm top}^2$, and $\Res$ has unique highest part $U_{\rm top}^2$. They cannot cancel. In the main norm the formal highest terms $a^2c^2$ cancel identically:
\[
N_1=2ac(X+\gamma H)+(X+\gamma H)^2-Hc^2.
\]
Its highest part comes from $2ac\gamma H$, of degree 18 or 35. The highest parts of $N_k$ and $N_t$ are nonzero. As the polynomial ring is an integral domain, the product degrees are exactly $18+7+2+488=515$ and $35+13+4+952=1004$.

For the integer form the positive-coordinate substitutions have degree two, $w=u^2$ has degree two, and the odd affine substitutions have degree one. Their disjoint leading forms preserve nonzero polynomials: setting the second and third argument of each $\Pos$ to zero leaves injective monomial substitutions $x\mapsto x_0^2$, followed by nonzero scalar substitutions for $i,t$. Thus the degree proof survives the complete substitution.

\subsection{Integer operation count}
The 27-integer circuit adds to the 84-gate positive circuit:
\begin{itemize}[nosep]
\item seven copies of $\Pos$, each with 7 gates: 49 gates;
\item the substitutions $i=2r_i+1$, $w=u^2$, $t=2r_t+1$: 5 gates;
\item the extra doubling $M=2M_0$: 1 gate.
\end{itemize}
The total is $84+49+5+1=139$, comprising 74 multiplications and 65 additions/subtractions. The integer system has 133 gates. The positive product has 49 multiplications and 35 additions/subtractions, and its system has 46 and 32, respectively.

\section{A smaller local compiler alphabet}
\label{sec:tiles}
The algebraic reductions above work with an authentic tuple independently of its height. We now reduce the tuple's defining local language. This changes neither the 215-state machine nor the arithmetic witness count.

\subsection{Four-cell tiles and the additional restrictions}
Retain package 11's 439 cell symbols. Ordinary symbols are $b\in\{0,1\}$ with no head, or $2(q+1)+b$ with a head in state $0\le q\le215$. The special symbols are
\[
L=434,\quad Q_s=435,\quad R_s=436,\quad V_b=437,\quad H_b=438.
\]
The phase symbols $L,Q_s,R_s$ have blank payload, with the initial head only at $Q_s$; $V_b,H_b$ are boundaries. A four-cell tile is $(l,c,r,n)$, where $n$ is the next-time center.

The inherited T4 predicate imposes at most one head in $(l,c,r)$, persistent vertical boundaries, horizontally uniform horizontal boundaries inside a strip, initial phase edges
\[
LL,\quad LQ_s,\quad Q_sR_s,\quad R_sR_s,
\]
no head escape through a vertical boundary, a halted head before a horizontal boundary, and the exact deterministic next-time update elsewhere. Initialization is tested on the outgoing edge from $H_b$. The executable predicate in the archive states every case; its inherited semantic proof is bundled~\cite{t4}.

We keep only tiles satisfying these further local rules:
\begin{enumerate}[nosep]
\item Each horizontal edge $(a,b)$ obeys: if $a=V_b$ then $b\in\{0,L,H_b\}$; if $b=V_b$ then $a\in\{0,R_s,H_b\}$. Away from $V_b$, an edge involving $H_b$ is $H_bH_b$; an edge involving a phase is one of the four displayed phase edges; otherwise both cells are ordinary.
\item If an ordinary center is adjacent to $V_b$, its other neighbor is headless.
\item If the next-time center is $H_b$, every head in $(l,c,r)$ is in state 215.
\end{enumerate}
These restrictions remove irrelevant locally admissible patterns while preserving sufficiently padded genuine computations.

\subsection{Soundness and completeness of the pruning}
\begin{proposition}
\label{prop:pruned}
A marked periodic array of the pruned tiles exists if and only if $K_{215}$ halts on blank tape. Its canonical completeness witness can have independently chosen sufficiently large width and height that are powers of five.
\end{proposition}
\begin{proof}
Soundness follows by inclusion in the inherited T4 language and its finite-strip extraction argument. This argument needs neither a globally unique Start nor coprimality of the cyclic stride and word length.

For completeness take any finite blank halting computation. Place at least two never-visited blank cells between every head position in the entire run and either vertical boundary. Extend the halted history using the absorbing completion $(215,b)\mapsto(b,0,215)$, which is used only after the original halting transition. Width and height may be increased independently.

In an interior row all horizontal edges stay within one type: initial phase, ordinary computation, or horizontal boundary. Edges at $V_b$ have the permitted blank, phase-margin, or horizontal-boundary endpoint. If an ordinary center is adjacent to $V_b$, the second margin cell is also permanently blank and headless. The row immediately before $H_b$ is a padded halted row, so every head in a local triple is halted.

Pass to the usual helical extension. An edge avoiding $V_b$ remains in a single row of the padded computation. At $V_b$, neighboring strips can differ by one time unit; the allowed boundary endpoint sets were chosen separately on the two sides and tolerate exactly this seam. The two permanently blank margin cells guarantee the neighboring-head restriction across every seam. Thus all added predicates hold. The original arbitrary padding and powers-of-five choice remain available.
\end{proof}

The three overlaps for cyclic tile positions are
\[
r_i=c_{i-1},\qquad c_i=l_{i-1},\qquad n_i=c_{i-h}.
\]
They give $w_i=(a_{i+1},a_i,a_{i-1},a_{i-h})$ for $a_i=c_i$, even when residues collide. Pulling back by $a_{-x-ht}$ gives the required plane. No extra arithmetic divisibility assumption is introduced.

\subsection{Distinguished tiles and exact census}
Put
\[
S=(434,435,436,1),\qquad E_*=(437,438,437,435).
\]
Both are locally valid. $S$ is the true initial neighborhood and its correct first successor. The tile $E_*$ is globally impossible: its two vertical neighbors persist to the next row, where its next center $Q_s$ would be adjacent to $V_b$ on both sides; the initial-phase rules would require it to be both $L$ and $R_s$. Thus forbidding $E_*$ removes no computation.

Order the tiles by $S$ first, $E_*$ second, then the remaining valid quadruples lexicographically. An exact enumeration and a separate class census give
\begin{center}
\begin{tabular}{@{}lr@{}}
\toprule
Center class & Number of tiles\\\midrule
Headless 0 & 1754\\
Headless 1 & 1744\\
430 working-head symbols, four each & 1720\\
Two halted-head symbols, eight each & 16\\
$L,Q_s,R_s$ & $6+1+6$\\
$V_b,H_b$ & $9+12$\\\midrule
Total & \textbf{5268}\\\bottomrule
\end{tabular}
\end{center}
The tile list is fully materialized in the archive, and its SHA-256 is
\begin{center}\small\code{6b53241de552096888a80c4d1947d061c32c1efac1f014ffedccfd06f5a6ea47}.\end{center}
Enumeration is finite, uses the literal transition table, and does not decide whether its blank computation halts.

\section{Exact coefficient recipe}
\label{sec:coeff}
The following definition makes the polynomial completely fixed, without materializing its enormous integer coefficients. All exponents here are fixed integers.

Set
\begin{gather}
a_* =439,\quad k_*=5268,\quad A_*=2^{13},\quad
\rho=2^{78125},\quad B=\rho^{9765625}=2^{5^{17}},\label{eq:baseradix}\\
M_*=11855,\quad E_{\max}=320085,\quad H_C=605923,\notag\\
T_1=1246533,\qquad T_2=1886707.\notag
\end{gather}
The four copy slots, in order, are $(1,0),(1,1),(1,2),(2,1)$; define
\[
e(r,c,u)=k_*+3a_*+(2-r)3a_*+(2-c)a_*+u,
\qquad0\le u<439.
\]
Thus their starting positions are 8780, 8341, 7902, and 7024. Let $\mathcal E$ contain selectors $0,\ldots,5267$, all four copy-slot positions, two dummy positions 10536 and 10537, and anchors $M_*,3M_*,9M_*,27M_*$. It has 7030 elements.

Put
\[
\mu=A_*-2+\sum_{j=1}^{2843}A_*^j.
\]
For tile $w_i=(w_i[0],\ldots,w_i[3])$, attach native coefficient
\[
v_i=1+\sum_{j=0}^{3}A_*^{1+ja_*+w_i[j]}
          +\sum_{t=1}^4A_*^{4a_*+t}.
\]
The copy coefficient at slot $j$, value $u$, is $A_*^{1+ja_*+u}$. The coefficient at anchor $3^{t-1}M_*$ is $A_*^{4a_*+t}$. Dummy coefficients are zero. Now define
\begin{align}
D_C={}&\rho^{3a_*}+\rho^{H_C+a_*}+\rho^{8M_*}+\rho^{24M_*}
 +\sum_{e\in\mathcal E}v_e(\rho^{T_1-e}+\rho^{T_2-e}),\notag\\
D_R={}&\rho^{H_C},\notag\\
M_F^0={}&\mu(\rho^{T_1}+\rho^{T_2})
       +\sum_{u=0}^{438}\rho^{e(1,1,u)}\notag\\
 &+\sum_{u=0}^{438}\left(\rho^{H_C+e(1,0,u)}+
                              \rho^{H_C+e(1,1,u)}\right)
       +\rho^{9M_*}+\rho^{27M_*},\notag\\
M_C={}&B-1-\sum_{\substack{e\in\mathcal E\\e\ne1}}\rho^e
                      -2\rho^{10536},\notag\\
M_F={}&M_F^0+4.\label{eq:coefficients}
\end{align}
The four ports are exactly
\begin{equation}
\boxed{(\beta,K_0,M_C,M_F^*)=(B-1,D_C+BD_R,M_C,M_F+B-1).}\label{eq:ports}
\end{equation}
Only the \emph{first} dummy gets the additional upper binary bit in $M_C$. Subtracting $2\rho^e$ at both dummies would break the population identity.

\subsection{Why the smaller layout satisfies the compiler contract}
The new layout is the same four-copy-slot construction as T4, with recomputed integers. Its parameters are
\[
\ell_*=13,\quad p_*=1083,\quad m_*=7029,\quad
b=5^7=78125,\quad L_*=5^{10}=9765625.
\]
We have $\pc(\mu)=12+2843=2855$, and the two clause bands, three copy-overlap groups, two tested anchors, and low marker give
\[
\pc(M_F)=2(2855)+3(439)+2+1=7030,
\quad\pc(M_C)=5^{17}-7030.
\]
This proves the population equality in \eqref{eq:masks}; the residues and strict inequalities there follow from the same disjoint bit positions.

For the inherited mass estimate put
\[
U_*=4(m_*+3)\bigl(18(k_*+4a_*+4)+6\bigr)+8.
\]
The selector coefficients have nine $A_*$-monomials and the other nonzero coefficients have one. The sufficient no-carry bound is
\[
b\ge\max\{13(2844)+3,\ 13(1760)+\operatorname{bitlength}(U_*)+1,\ 4\}
=36975,
\]
which $b=78125$ satisfies. The clause-band difference is
\[
T_2-T_1=640174>2E_{\max}=640170,
\quad T_2-T_1\equiv2\pmod4.
\]
All bands and anchor translations fit inside $L_*$. The archive checks the actual support, all 1319 copy/anchor field targets, separation of both clause bands, and unique anchor alignment at radix positions. The separate parity argument in the compiler proof excludes nonaligned binary rotations; the program does not enumerate the enormous set of all bit rotations. These are the same hypotheses used in the inherited native compiler proof; the support and mass certificates establish their finite parts for this tuple. Also $2D_C-D_R\equiv3\pmod5$, so the five-adic dummy tuning remains valid.

\subsection{Canonical parity survives the compiler change}
\label{sec:parity}
Choose both spatial width and temporal height as sufficiently large powers of five, which Proposition~\ref{prop:pruned} permits. Their product $N_*$ is odd, and the bit length exponent of the canonical packed base is $t_*=5^{17}N_*$, also odd. The canonical packed word has $C_w\equiv1\pmod4$, and $M_C\equiv2\pmod4$. At that witness, $q$ is divisible by four, $J\equiv1\pmod4$, and the rank formula gives
\[
R\equiv C_w+M_C\equiv3\pmod4.
\]
The canonical relation $X=2^R$ and $q=2^{t_*}$ therefore gives
\[
w=2^{R-t_*},
\]
an integer square because $R-t_*$ is positive and even.

The canonical $A_0$ is even and $R$ is odd, hence $c=\psi_{A_0}(R)$ is odd. Use strong index $m=cR$. The binomial expansion of
\[
(\chi_{A_0}(R)+c\sqrt\Delta)^c
\]
shows $c^2\mid\psi_{A_0}(cR)$: its first square-root term contains $c\cdot c$, and every later term contains at least $c^3$. Thus
\[
i=\frac{\psi_{A_0}(cR)}{c^2}
\]
is an odd positive integer. The canonical auxiliary root is $y=\psi_{\rho_0}(R)$ for its auxiliary Pell base $\rho_0$. The recurrence modulo two gives $\psi_{\rho_0}(n)\equiv n\pmod2$, so $y$ is odd. These are exactly the completeness properties required for \eqref{eq:integerweight}.

\section{Verification and reproducibility}
\label{sec:checks}
The archive is self-contained for its finite checks. The construction programs use Python's standard library and small pinned inputs. They never expand $B$, search for a refutation of ZFC, or simulate an unknown halting run to completion.

The verification has distinct scopes:
\begin{enumerate}
\item \textbf{Circuit syntax and dimensions.} Every operand is defined, every register is unique, every arithmetic gate and witness port is live, and the operation counts match the emitted JSON programs.
\item \textbf{Independent algebra.} Direct formulas separate from the circuit evaluate the signed polynomial and its residuals on 100 exact integer assignments, including negative inputs. There are also 216 exact rational reconstruction fixtures with signed roots. These test algebra and signs, not the existence of a full K215 witness.
\item \textbf{Degree corroboration.} Modular polynomial restrictions corroborate the exact degree calculations. The uniform leading-form proof in Section~\ref{sec:degree}, rather than these finite restrictions, establishes the theorem for the actual coefficients.
\item \textbf{Literal machine and local language.} The compiler rechecks the table hash, all 430 actions, the inherited 37,075-tile count, and the new 5,268-tile list. Padded control computations test finite seams. Controls deliberately have known halting behavior and are not presented as runs of the unmodified K215.
\item \textbf{Coefficient geometry.} Native support, mass bounds, field targets, population identities, residues, anchor synchronization, and the fixed coefficient binding are checked without expanding the final numerals.
\end{enumerate}
The mathematical proof of arbitrary-length soundness and completeness remains separate from finite checks. The existing Lean/Rocq endpoint was inspected through its pinned sources and verification records; a fresh endpoint build was not run here.

Run \code{python3 verify.py} from the extracted directory to regenerate the circuits and finite receipts. The checked results are saved under \code{verification/}. The article is rebuilt with \code{make pdf} when a standard LaTeX installation is available. \code{SOURCE_PINS.json} records exact upstream sources; \code{MANIFEST.sha256} records the delivered files. The directory can be incorporated as a new report package without changing the older packages.

\section{What remains open}
\label{sec:questions}
\begin{enumerate}
\item \textbf{Fewer than 27 unrestricted integers.} Seven coordinates still require explicit positivity. Removing a positive coordinate costs potentially three integer parameters, but its sign cannot simply be dropped: the transport and Pell interval arguments use these signs. A useful next target is a sound alternative to the two independent positive interval slacks $\eta,\zeta$.
\item \textbf{Fewer than 13 positive integers.} The present weight only merges three square roots into one coordinate. Further elimination must prevent all vanishing denominators and preserve integer, rather than merely algebraic, roots. An ordinary-strong version could lower degree, but clearing its $\Delta$ denominator requires a new integrality argument.
\item \textbf{A smaller degree at 27 integers.} The dominant term is the iterated root-resultant term involving $Qz_0^2$. Keeping one extra root may yield a much lower-degree 28-integer option. The present count improvement should not hide that tradeoff.
\item \textbf{Fewer than 67 evaluator operations.} This remains open for the collection's fixed-blank polynomial. Our positive evaluator keeps 84 gates while reducing degree. A separate circuit optimization objective is needed to improve the 67-gate baseline.
\item \textbf{Smaller coefficient geometry.} The new tile count decreases by a factor exceeding seven, but the cell length stays $5^{10}$; the radix improvement comes from lowering the inner bit radix from $5^8$ to $5^7$. Tighter anchors, band locations, or symbol coding could cross the next power-of-five threshold. Such a change must retain the all-rotation argument and exact mask populations.
\item \textbf{An explicit nine-unknown specialization.} The existing abstract nine-unknown representation theorem does not select and audit the concrete small equation sought here. Extracting a fully fixed specialization with a transparent coefficient recipe remains a distinct research objective.
\item \textbf{Formalizing the new bridge.} A small formal development can first prove the rational-square lemma, the weighted resultant inverse, the $T=0$ exclusion, and the parity substitution. A full proof-assistant result additionally needs the native compiler and its Pell interfaces; reusing the existing machine endpoint alone does not discharge them.
\end{enumerate}

\section{Conclusion of the construction}
The main deliverable is the fixed polynomial $P_{27}$ defined by \eqref{eq:common}--\eqref{eq:polynomial}, \eqref{eq:pos}, and \eqref{eq:integerweight}, with the single coefficient tuple \eqref{eq:ports}. Its improvement is concrete: 27 unrestricted integer variables instead of the pinned collection's 29. The same argument supplies a degree-515, 84-operation polynomial in 13 positive integers, and both use a smaller coefficient radix $2^{5^{17}}$.

The biconditionals state that these equations encode syntactic inconsistency. They do not prove that the equations have no solutions, and they do not prove $\Con(\mathrm{ZFC})$. Under the usual second-incompleteness hypotheses, consistency of ZFC rules out a ZFC proof of their unsolvability. Two-sided undecidability should not be asserted from consistency alone.

\begin{thebibliography}{12}\small\raggedright
\bibitem{collection} ProveIt contributors. \emph{Con(ZFC) in Diophantine form: wave-2 reconciliation}. Pinned commit \code{f0ae4e0991}. \href{https://github.com/VladimirReshetnikov/ProveIt/blob/f0ae4e0991ca44a7bb4e5671d42e33d775e48035/SetTheory/Cardinals/docs/reports/hilbert-tenth-problem/con-zfc-diophantine-forms/README.md}{Repository source}. Bundled comparison note.
\bibitem{amalgamation} Package 09-combined13. \emph{Root amalgamation for the fixed blank K215 slice}; arithmetic circuits and degree audits. \href{https://github.com/VladimirReshetnikov/ProveIt/blob/f0ae4e0991ca44a7bb4e5671d42e33d775e48035/SetTheory/Cardinals/docs/reports/hilbert-tenth-problem/con-zfc-diophantine-forms/09-combined13/proofs/ROOT_AMALGAMATION_PROOF.md}{Pinned proof}. Bundled in \code{upstream/}.
\bibitem{integer29} Package 14-integer29. \emph{Independent audit of signed-domain reduction}. \href{https://github.com/VladimirReshetnikov/ProveIt/blob/f0ae4e0991ca44a7bb4e5671d42e33d775e48035/SetTheory/Cardinals/docs/reports/hilbert-tenth-problem/con-zfc-diophantine-forms/14-integer29/proofs/signed_domain_proof.md}{Pinned proof}. Bundled in \code{upstream/}.
\bibitem{t4} Package 11-t4. \emph{Audit and smaller coefficient backend for the combined consistency equation}. \href{https://github.com/VladimirReshetnikov/ProveIt/blob/f0ae4e0991ca44a7bb4e5671d42e33d775e48035/SetTheory/Cardinals/docs/reports/hilbert-tenth-problem/con-zfc-diophantine-forms/11-t4/bridge/BRIDGE_AUDIT.md}{Pinned T4 proof}. Bundled in \code{upstream/}.
\bibitem{blank} Package 06-blank67, preserved in package 09. \emph{Independent proof and source audit: the 67-operation blank-input polynomial}; \emph{Blank helical proof}. Bundled as \code{FIXED_BLANK_PROOF.md} and \code{BLANK_HELICAL_PROOF.md}.
\bibitem{pell} Package 06-blank67, preserved in package 09. \emph{Pell lemma audit}. Bundled as \code{upstream/PELL_LEMMA_AUDIT.md}; its source dependencies are identified within the note.
\bibitem{signedquotient} ProveIt native-stream-queue research. \emph{Signed auxiliary quotient absorption and soundness}. \href{https://github.com/VladimirReshetnikov/ProveIt/tree/f0ae4e0991ca44a7bb4e5671d42e33d775e48035/Computability/HilbertTenthProblem/Papers/research-wip/native-stream-queue}{Pinned source directory}. Two pertinent proof notes are bundled.
\bibitem{machine} ProveIt contributors. \emph{K215 ordinary calibration proof}. \href{https://github.com/VladimirReshetnikov/ProveIt/blob/f0ae4e0991ca44a7bb4e5671d42e33d775e48035/Computability/BusyBeaver/Research/ZFC/candidate-k215/proof.md}{Pinned proof}. Bundled as \code{upstream/K215_CALIBRATION.md}.
\bibitem{endpoint} ProveIt contributors. \emph{KEndpoint215.lean}. \href{https://github.com/VladimirReshetnikov/ProveIt/blob/f0ae4e0991ca44a7bb4e5671d42e33d775e48035/Computability/BusyBeaver/Lean/BusyBeaver/ZFC/KEndpoint215.lean}{Pinned endpoint}. Bundled unmodified.
\bibitem{verification} ProveIt contributors. \emph{K215 verification records}. \href{https://github.com/VladimirReshetnikov/ProveIt/blob/f0ae4e0991ca44a7bb4e5671d42e33d775e48035/Computability/BusyBeaver/Research/ZFC/k225-formal/verification/k215/README.md}{Pinned scope and records}. Bundled as \code{upstream/K215_VERIFICATION.md}.
\end{thebibliography}
\end{document}
